В условиях стремительного развития технологий и возрастающих требований к эффективности \cite{grushezkii_2021, grushezkii_and_2015, repin_2008}, автоматизация рабочих процессов наземных транспортно-технологических машин (НТТМ) становится ключевым фактором повышения производительности и снижения затрат. Современные подходы к~автоматизации охватывают широкий спектр решений, начиная от внедрения роботизированных комплексов и заканчивая использованием интеллектуальных систем управления во всех группах машин (рис.\ \cref{fig:cls}). Роботизированная наземная транспортно-технологическая машина (РНТТМ) в настоящей работе рассматривается как, например, частный случай высокоавтоматизированной землеройной и дорожно-строительной машины (ВАЗДСМ) по \cite{gostR71615}, функционирующей на уровнях автоматизации УА2-УА5.

\nomenclature{НТТМ}{наземные транспортно-технологические машины\nomrefpage}
\nomenclature{РНТТМ}{роботизированные НТТМ\nomrefpage}
\nomenclature{ВАЗДСМ}{высокоавтоматизированная землеройная и дорожно-строительная машина\nomrefpage}

\begin{figure}[ht]
	\centering{
		\includegraphics[scale=0.5]{drawio/cls.drawio.png}
	}
	\caption{Классификация укрупненных групп НТТМ по технологическому назначению}\label{fig:cls}
\end{figure}

Например, в проектах строительства скоростных дорог М-11 и М-12 применяются следующие основные группы машин: экскаваторы (ЭО-4321, ЕК-270), бульдозеры (Б10ПМ, ДЗ-101А), автогрейдеры (ДЗ-98), катки грунтовые и асфальтобетонные (Bomag, Hamm), асфальтоукладчики (Vogele, Titan) и т.д.

Рабочий процесс отдельной НТТМ -- это совокупность взаимосвязанных операций, выполняемых машиной при непосредственном взаимодействии с рабочей средой (грунт, груз, дорожное покрытие, снег и др.) и опорной поверхностью, направленных на преобразование, перемещение или обработку объекта труда. К рабочим процессам НТТМ относятся:

\begin{itemize}
	\item транспортирование грузов;
	\item разработка, перемещение и уплотнение грунта;
	\item уборка территорий (снег, мусор);
	\item планировка и профилирование поверхностей;
	\item погрузочно-разгрузочные операции;
	\item специальные технологические операции (виброуплотнение, фрезерование и др.).
\end{itemize}

Рассмотрим такой элемент работ в строительном производстве, как транспортирование грузов в технологической цепочке доставка-переработка-погрузочно-разгрузочные работы дорожно-строительных материалов и т.\ д. Он занимает около 42\% всей трудоемкости дорожно-строительных работ \cite{Vilman, grushezkii_and_2021}. Рациональное совершенствование транспортных и~погрузочных операций является одним из важнейших направлений по сокращению затрат ручного труда в строительстве.

Так, наиболее близким примером автоматизации к сфере НТТМ в этой области является высокоавтоматизированное транспортное средство (ВТС).

\nomenclature{ВТС}{высокоавтоматизированное транспортные средства\nomrefpage}

В мире существует большое количество компаний с мировым именем, чья деятельность так или иначе связана с разработкой ВТС: Renault, Audi (Volkswagen Group), Baidu, BMW, Intel и Mobileye, Toyota, Uber, DAF, Daimler, Iveco, MAN, Ford,  Google, Fiat, Chrysler Automobiles, Honda, Hyundai, Jaguar, Land Rover, General Motors, Lyft, Bosch, Mercedes-Benz (Daimler AG), Apple, Microsoft, Nissan, Scania, Volvo, Nvidia, PSA Groupe, Tata Elixsi, Tesla, Volkswagen, Volvo, Yutong и~др.

Российское правительство и промышленность сегодня также обращают повышенное внимание к концепции создания ВТС: КамАЗ, Яндекс, ФГУП НАМИ и~Кировский завод, Волгабас, ведут проекты по созданию ВТС в~партнерстве с разработчиками систем автоматического управления (Google, Cognitive Technologies, Uber, Baidu, КБ Аврора и другие) \cite{AleshinMendozaTarasovKlyavin2018}.

Анализ показывает \cite{Saikin}, что разработка ВТС происходит по двум основным направлениям:

\begin{itemize}
    \item внедрение и расширение функциональности различных систем помощи водителю (которыми в настоящее время серийно комплектуются автомобили всех классов);

    \item создание способов и систем управления полностью ВТС (которые в~настоящее время находятся на стадии испытаний опытных образцов, в~том числе эксплуатационных).
\end{itemize}

На рисунке~\cref{fig:patent} представлена динамика разработок наземных ВТС, усовершенствованных систем помощи водителю (ADAS) и их компонентов за рубежом.

\nomenclature{ADAS}{Advanced Driver-Assistance Systems, усовершенствованная система помощи водителю\nomrefpage}

\begin{figure}[ht]
    \centering{
        \includegraphics[scale=0.4]{views/patent.pdf}
    }
    \caption{Динамика разработок за рубежом наземных ВТС, систем ADAS и~их компонентов}\label{fig:patent}
\end{figure}

Для реализации автономного управления движением наиболее активно ведутся разработки и патентуются системы управления движением, которые включают: блоки управления интеллектуальным шасси, системы технического зрения, обработки и передачи информации, системы навигации и ориентации (рис.\ \cref{fig:system}).

\begin{figure}[ht]
    \centering{
        \includegraphics[scale=0.5]{drawio/system.drawio.png}
    }
    \caption{Подсистемы автоматизированной системы управления ВТС}\label{fig:system}
\end{figure}

Конкретным примером в сфере РНТТМ является разработка Научно-технического центра ПАО <<КАМАЗ>> совместно со специалистами МГТУ им.~Н.~Э.~Баумана: транспортное средство КАМАЗ-6559 (рис.\ \cref{fig:ypiter30}), предназначенный для перевозки сыпучих грузов в автономном режиме.

\begin{figure}[ht]
    \centering{
        \includegraphics[scale=0.35]{views/kamaz.png}
    }
    \caption{Беспилотный КАМАЗ-6559 (<<Юпитер 30>>)}\label{fig:ypiter30}
\end{figure}

ВТС создан в рамках проекта «Создание семейства электромеханических беспилотных автомобилей-самосвалов большой грузоподъемности в интересах добывающих отраслей промышленности РФ». Предназначен для перевозки разрыхленной горной массы или руды по безлюдной технологии -- без присутствия людей в опасной зоне работы большегрузных машин и экскаваторов. В связи с~этой особенностью в конструкции автомобиля отсутствует кабина для водителя. РНТТМ изначально проектировался как автономное транспортное средство для работы в карьере. Имеет специальные, защищенные от пыли, грязи, влаги и вибраций видеокамеры, 2D- и 3D-лидары, ультразвуковые датчики, радары. Также имеются GSM-антенны, система навигации Global Positioning System (GPS) и Глобальная навигационная спутниковая система (ГЛОНАСС).

\nomenclature{GSM}{Global System for Mobile Communications\nomrefpage}
\nomenclature{GPS}{Global Positioning System\nomrefpage}
\nomenclature{ГЛОНАСС}{Глобальная навигационная спутниковая система\nomrefpage}

Другим примером РНТТМ является наземный робототехнический комплекс <<Курьер>> производства компании <<Садаев Групп>> (рис.\ \cref{fig:kurier}), который выполняет специальные задачи. Например: эвакуацию раненых с поля боя или доставку различных грузов в труднодоступные места.

\nomenclature{НРТК}{наземный робототехнический комплекс\nomrefpage}

\begin{figure}[!h]
	\centering{
		\includegraphics[scale=4]{views/kurier.png}
	}
	\caption{Наземный робототехнический комплекс <<Курьер>>}\label{fig:kurier}
\end{figure}

Еще одним примером является автономный робот <<Пиксель>> (рис.\ \cref{fig:pixel1}, \cref{fig:pixel2}) разработки российской компанией <<Автономика>>, которые c 2024 года участвуют в регулярном обслуживании дорог в городах Москва \footnote{\ \url{https://www.mos.ru/mayor/themes/10719050/}} и Санкт-Петербург \footnote{\ \url{https://www.tadviser.ru/index.php/Продукт:Пиксель\_(робот-уборщик)}}.

\begin{figure}[!h]
	\centering{
		\includegraphics[scale=0.3]{views/pixel1_1.png}
	}
	\caption{РНТТМ <<Пиксель>> с щеткой}\label{fig:pixel1}
\end{figure}

\begin{figure}[!h]
	\centering{
		\includegraphics[scale=0.3]{views/pixel1_2.png}
	}
	\caption{РНТТМ <<Пиксель>> со скребком}\label{fig:pixel2}
\end{figure}

При определении подходов к автоматизации рабочих процессов РНТТМ необходимо учесть их особенности от обычных автономных транспортных средств в части, касающихся:

\begin{itemize}
\item \textit{размеров}: совершенно иные габариты и форма кузова налагают требования на конструкции элементов автоматизированной системы управления РНТТМ;

\item \textit{управляемости}: шасси РНТТМ существенно отличается от других ТС с точки зрения управления перемещением: массой, динамических характеристик, гораздо большие углы (крены) рамы и кабины;

\item \textit{условий эксплуатации}: для РНТТМ параметры иные: ТС классов 1 или 2 могут передвигаться со скоростью не менее 70 км/ч по загородным дорогам и 60 км/ч в городском режиме, движение РНТТМ зависит от его типа и конструкции (колесный или гусеничный и т.\ д.). Изменения рабочей скорости напрямую влияет на требования к сенсорам: какая должна быть дальность обнаружения объектов, частота опроса датчиков и алгоритмы детекции;

\item \textit{компоновки вычислительных средств и систем}: она также сильно зависит от типа кузова, более высоких требований к защите от внешних воздействий (у РНТТМ гораздо жестче вибрации, возможны ударные воздействия), наличия сменных компонентов в структуре оборудования, а~также требуется резервный контур в системе управления;

\item \textit{дополнительных конструкций}: например наличие полуприцепа (другого рабочего оборудования), что в несколько раз усложняет алгоритмы управления движением: существенно меняется модель поведения, ограничиваются обзоры видимости для сенсоров. Он имеет свою массу, которая может меняться в диапазоне нескольких тонн.
\end{itemize}